\begin{helper}[A distinct tree-meeting two-color component is a cycle]
Let (T) be a tree-sequence in a finite loopless multigraph and let
\(\varphi\) be a proper partial edge coloring.  Fix two colors
\(\alpha,\beta\) and two nodes (z,w) of their two-color incidence graph.
Suppose the component through (z) is a path component, the two colors are
interchangeable on (T), both components meet (V(T)), and the two
components are distinct.  Then the component through (w) is a two-color
cycle: there is a simple closed walk whose support is exactly that component.
\end{helper}
\begin{proof}
By \noderef{TwoColorPathComponent}, the path-component hypothesis includes
(\alpha\ne\beta).  If the component through (w) were also a path component,
the definition of
\noderef{TreeInterchangeable}, applied to the two supplied tree-intersection
witnesses, would identify it with the component through (z), contrary to
the distinctness hypothesis.  Hence the component through (w) has no node
whose neighbor set is a subsingleton.

By \noderef{TwoColorComponentSpanningPath}, properness gives a simple path
(R) whose support is the whole component through (w).  This path is not
trivial: otherwise its sole node would have an empty neighbor set inside its
connected component and hence, since the support is the whole component, in
the incidence graph as well.  That neighbor set would be a subsingleton,
contrary to the preceding paragraph.

Let (a,b) be the two endpoints of (R), and let (a_1) be the path neighbor of
(a).  The neighbor set of (a) is not a subsingleton, so it has a neighbor
(a') distinct from (a_1).  Every neighbor remains in the same connected
component and hence occurs on (R).  The node (a') cannot be internal to
(R): an internal node already has its two distinct path neighbors, and its
additional neighbor (a) would give it at least three neighbors, contrary to
\noderef{TwoColorIncidenceDegreeAtMostTwo}.  Nor can (a') equal (a), since
the incidence graph is loopless.  The only remaining possibility is the
other endpoint (b).  In particular (b\ne a_1), so the path has more than one
edge, and (a) and (b) are adjacent by an incidence edge not already used by
(R).

Close (R) with this last adjacency.  Simplicity of (R), together with the
fact that its endpoints are distinct and the closing incidence edge is not a
path edge, makes the resulting closed walk a simple cycle.  Closing the path
does not add a node, so its support is still the support of (R), namely the
entire component through (w).  This is exactly
\noderef{TwoColorComponentCycle}.
\end{proof}
